% ===== Block B04: 原文 p.31–36 | §1.4 稳定带点曲线的通用族 | 译手 tr-ch1c =====
\subsection{稳定带点曲线的通用族}
% [原文 p.31]

让我们考虑亏格为 $g$ 的稳定 $N$-点曲线的一个通用族
$\mathfrak{F}=(\pi:\mathcal{C}\to\mathcal{B};s_1,s_2,\dots,s_N)$。我们假定已有以
$s_N(\mathcal{B})$ 为中心的一个局部（形式）坐标 $\eta$。下文我们需要局部地考虑族
$\mathfrak{F}$。为此，我们引入 $\mathcal{C}$ 的如下局部坐标。

\begin{sloppypar}对 $\pi$ 的临界轨迹 $\Sigma$ 上的点 $P$，我们可以选取 $\mathcal{C}$ 的以 $P$
为中心的局部坐标 $(u_1,u_2,\dots,u_{M-1},z,w)$，以及 $\mathcal{C}$ 的以 $\pi(P)$
为中心的局部坐标 $(\tau_1,\tau_2,\dots,\tau_M)$，使得全纯映射 $\pi$ 由
\end{sloppypar}
\[
(u_1,u_2,\dots,u_{M-1},z,w)\longmapsto
(u_1,u_2,\dots,u_{M-1},zw)=(\tau_1,\tau_2,\dots,\tau_M)
\]
给出。换言之，我们有
\[
\pi^*\tau_k=
\begin{cases}
u_k, & k=1,2,\dots,M-1,\\
zw, & k=M.
\end{cases}
\]

对点 $P\in\mathcal{C}\setminus\Sigma$，我们可以选取 $\mathcal{C}$ 的以 $P$
为中心的局部坐标 $(u_1,u_2,\dots,u_M,z)$，以及 $\mathcal{C}$ 的以 $\pi(P)$
为中心的局部坐标 $(\tau_1,\tau_2,\dots,\tau_M)$，使得全纯映射由
\[
(u_1,u_2,\dots,u_{M-1},z)\longmapsto
(u_1,u_2,\dots,u_M)=(\tau_1,\tau_2,\dots,\tau_M)
\]
给出。
% [Tutor Note: 这里两处原文都写 "local coordinates … of $\mathcal{C}$ with center
%   $\pi(P)$"，但 $\pi(P)\in\mathcal{B}$，应为 "of $\mathcal{B}$"；又第二个显示式源空间
%   多写了一个坐标 $z$（$\mathcal{C}$ 维数为 $M+1$，$\mathcal{B}$ 维数为 $M$），按
%   $\pi$ 的局部同构性应为 $(u_1,\dots,u_M,z)\mapsto(u_1,\dots,u_M)$。两处均照原文排，
%   见报告 ERRATA。]

$\mathcal{O}_{\mathcal{C}}$-模 $\Omega^1_{\mathcal{C}/\mathcal{B}}$ 由正合序列
\[
\pi^{-1}\Omega^1_{\mathcal{B}}\otimes_{\mathcal{O}_{\mathcal{B}}}
\mathcal{O}_{\mathcal{C}}\to\Omega^1_{\mathcal{C}}\to
\Omega^1_{\mathcal{C}/\mathcal{B}}\to 0
\]
定义。
% [Tutor Note: 这是 $\Omega^1_{\mathcal{C}/\mathcal{B}}$ 的定义式（相对微分模），
%   等价说法：它是 $\Omega^1_{\mathcal{C}}$ 模掉由 $\pi^{-1}\Omega^1_{\mathcal{B}}$
%   的像生成的子模。核对维数：在 $P\notin\Sigma$ 处 $\Omega^1_{\mathcal{C}/\mathcal{B}}$
%   秩为 $1$，在结点处也秩为 $1$（由下面 (1.57) 可见），故它总是局部自由的。]

层 $\Omega^1_{\mathcal{C}/\mathcal{B}}$ 称为族 $\pi:\mathcal{C}\to\mathcal{B}$
的相对 $1$-形式的芽层。让我们利用上述局部坐标来描述层
$\Omega^1_{\mathcal{C}/\mathcal{B}}$。在点 $P\in\mathcal{C}\setminus\Sigma$
的邻域内，层 $\Omega^1_{\mathcal{C}/\mathcal{B}}$ 局部同构于
$\mathcal{O}_{\mathcal{C}}\dd z$。在 $P\in\Sigma$ 的小邻域内，我们有
$\mathcal{O}_{\mathcal{C}}$-模同构
\begin{equation}
\Omega^1_{\mathcal{C}/\mathcal{B}}\simeq
(\mathcal{O}_{\mathcal{C}}\dd z+\mathcal{O}_{\mathcal{C}}\dd w)
\big/\mathcal{O}_{\mathcal{C}}(w\,\dd z+z\,\dd w). \tag{1.57}
\end{equation}
% [Tutor Note: (1.57) 的核对办法：在结点处 $\mathcal{O}_{\mathcal{C}}$ 的局部模型是
%   $\C\{z,w\}/(zw)$，而 $\dd(zw)=w\,\dd z+z\,\dd w$ 正是被模掉的子模，它对应"相对"
%   微分（即沿纤维方向的微分）。在 $z\ne0$ 的枝上 $\dd z$ 生成该模，在 $w\ne0$ 的枝上
%   $\dd w$ 生成，因为此时 $\dd w=-(w/z)\dd z$。这个"两枝粘接"的看法在 p.32 的
%   $\omega_{\mathcal{C}/\mathcal{B}}$ 描述里再次出现。]

此外，我们有下述引理。

\begin{lemma}[1.31]
下述序列
\begin{equation}
0\to\pi^{-1}\Omega^1_{\mathcal{B}}\otimes_{\mathcal{O}_{\mathcal{B}}}
\mathcal{O}_{\mathcal{C}}\to\Omega^1_{\mathcal{C}}\to
\Omega^1_{\mathcal{C}/\mathcal{B}}\to 0 \tag{1.58}
\end{equation}
是正合的，并且给出层 $\Omega^1_{\mathcal{C}/\mathcal{B}}$ 的一个局部自由分解。
\end{lemma}

设 $\omega_{\mathcal{C}/\mathcal{B}}$ 为 $\pi:\mathcal{C}\to\mathcal{B}$ 的
相对对偶层。由于 $\mathcal{C}$ 与 $\mathcal{B}$ 非奇异且 $\pi$ 平坦，我们有
$\mathcal{O}_{\mathcal{C}}$-模同构
\[
\omega_{\mathcal{C}/\mathcal{B}}\simeq
\omega_{\mathcal{C}}\otimes(\pi^*\omega_{\mathcal{B}}^{-1}),
\]
其中 $\omega_Y$ 是复流形 $Y$ 的对偶层（典范层）。相对对偶层
$\omega_{\mathcal{C}/\mathcal{B}}$ 局部描述如下。（亦见 (1.44)。）

在点 $P\in\mathcal{C}\setminus\Sigma$ 的小邻域内，我们有
\[
\omega_{\mathcal{C}/\mathcal{B}}=
\Omega^1_{\mathcal{C}/\mathcal{B}}\simeq\mathcal{O}_{\mathcal{C}}\dd z.
\]
在点 $P\in\Sigma$ 的小邻域内，我们有
\[
\omega_{\mathcal{C}/\mathcal{B}}\simeq
\mathcal{O}_{\mathcal{C}}(\dd z\wedge\dd w)\otimes(\dd\tau_M)^{-1}.
\]
% [Tutor Note: 在光滑点，相对对偶层就是相对微分层，故 $\omega_{\mathcal{C}/\mathcal{B}}
%   =\Omega^1_{\mathcal{C}/\mathcal{B}}\simeq\mathcal{O}_{\mathcal{C}}\dd z$。在结点处
%   这个公式要按"形式分式"理解：$\dd\tau_M=\dd(zw)=z\,\dd w+w\,\dd z$，而
%   $(\dd z\wedge\dd w)\otimes(\dd\tau_M)^{-1}$ 正是 p.32 用来定义 $\omega$ 的生成元的
%   写法。$\dd z\wedge\dd w$ 被 $(z,w)$ 零化（因为 $z\,\dd z\wedge\dd w=0$
%   = $w\,\dd z\wedge\dd w$），这一点是核对下面粘接关系的关键。]

% [原文 p.32]

特别地，我们有
\[
\omega_{\mathcal{C}/\mathcal{B}}\simeq
\begin{cases}
\mathcal{O}_{\mathcal{C}}\dfrac{\dd z}{z}, & \text{若 } z\ne 0,\\[6pt]
\mathcal{O}_{\mathcal{C}}\dfrac{\dd w}{w}, & \text{若 } w\ne 0,
\end{cases}
\]
并带有关系
\[
\frac{\dd z}{z}+\frac{\dd w}{w}=0
\qquad\text{若 } zw\ne 0.
\]
% [Tutor Note: 这就是结点曲线对偶层的经典描述：在结点的两枝上分别取留数为 $+1$ 与
%   $-1$ 的简单极点微分。核对：由 $\dd(zw)=0$ 得 $w\,\dd z=-z\,\dd w$，两边除以 $zw$
%   即得 $\dd z/z=-\dd w/w$，也就是上面那条关系。也说明 $\omega_{\mathcal{C}/\mathcal{B}}$
%   在结点处是秩 $1$ 自由的（生成元 $\dd z/z$），而 $\Omega^1_{\mathcal{C}/\mathcal{B}}$
%   在结点处不是自由的（(1.57) 的分母不是主理想）。]

\begin{lemma}[1.32]
存在正合序列
\[
0\to\Omega^1_{\mathcal{C}/\mathcal{B}}\to\omega_{\mathcal{C}/\mathcal{B}}
\to\omega_{\mathcal{C}/\mathcal{B}}\otimes_{\mathcal{O}_{\mathcal{C}}}
\mathcal{O}_{\Sigma}\to 0.
\]
\end{lemma}

\noindent\textit{证明.}\quad 该映射 $\Omega^1_{\mathcal{C}/\mathcal{B}}\to
\omega_{\mathcal{C}/\mathcal{B}}$ 在点 $P\in\mathcal{C}\setminus\Sigma$ 的
邻域内局部由
\[
\dd z\longmapsto\dd z
\]
给出；而在点 $P\in\Sigma$ 的邻域内由
\[
\dd z\longmapsto z(\dd z\wedge\dd w)\otimes(\dd\tau_M)^{-1},
\]
\[
\dd w\longmapsto w(\dd z\wedge\dd w)\otimes(\dd\tau_M)^{-1}
\]
给出。
% [Tutor Note: 关键核对：把 (1.57) 的分母 $w\,\dd z+z\,\dd w$ 代入上面两式，得
%   $w\cdot z(\dd z\wedge\dd w)\otimes(\dd\tau_M)^{-1}
%   +z\cdot w(\dd z\wedge\dd w)\otimes(\dd\tau_M)^{-1}
%   =2zw(\dd z\wedge\dd w)\otimes(\dd\tau_M)^{-1}$。它看似不为零，但在
%   $\omega_{\mathcal{C}/\mathcal{B}}$ 中 $\dd z\wedge\dd w$ 被 $(z,w)$ 零化是在
%   $\Omega^2$ 层里成立的关系，此处应改用 p.32 的粘接描述来核验：在 $z\ne0$ 的枝上
%   像为 $z\cdot\dd z/z=\dd z$，与光滑情形一致；在 $w\ne0$ 的枝上像为
%   $w\cdot\dd w/w=\dd w$，故该映射无混淆。]

特别地，我们有
\[
\dd z\longmapsto z\,\frac{\dd z}{z}\qquad\text{若 } z\ne 0,
\]
\[
\dd w\longmapsto w\,\frac{\dd w}{w}\qquad\text{若 } w\ne 0.
\]
这就证明了引理 1.32。\hfill$\square$

\begin{lemma}[1.33]
令
\begin{equation}
\Theta_{\mathcal{C}/\mathcal{B}}=
\underline{\mathrm{Hom}}_{\mathcal{O}_{\mathcal{C}}}
(\Omega^1_{\mathcal{C}/\mathcal{B}},\mathcal{O}_{\mathcal{C}}). \tag{1.59}
\end{equation}
则 $\Theta_{\mathcal{C}/\mathcal{B}}$ 是一个可逆
$\mathcal{O}_{\mathcal{C}}$-模，并且存在同构
\begin{equation}
\Theta_{\mathcal{C}/\mathcal{B}}\simeq
\underline{\mathrm{Hom}}_{\mathcal{O}_{\mathcal{C}}}
(\omega_{\mathcal{C}/\mathcal{B}},\mathcal{O}_{\mathcal{C}}). \tag{1.60}
\end{equation}
因此，$\Theta_{\mathcal{C}/\mathcal{B}}$ 是可逆层。
\end{lemma}

\noindent\textit{证明.}\quad 由 (1.58) 容易看出，在点 $P\in\mathcal{C}\setminus\Sigma$ 的
邻域内我们有同构
\[
\Theta_{\mathcal{C}/\mathcal{B}}\simeq
\mathcal{O}_{\mathcal{C}}\frac{\partial}{\partial z},
\]
而在点 $P\in\Sigma$ 的邻域内我们有同构
\begin{equation}
\Theta_{\mathcal{C}/\mathcal{B}}\simeq
\mathcal{O}_{\mathcal{C}}\left(z\frac{\partial}{\partial z}
-w\frac{\partial}{\partial w}\right). \tag{1.61}
\end{equation}
由这一事实与 (1.57)，我们得到所要的结果。\hfill$\square$
% [Tutor Note: (1.61) 的符号可以这样核对：相对向量场必须保持纤维方程
%   $zw=\pi^*\tau_M$，即 $v(zw)=0$（模去从 $\mathcal{B}$ 拉回的项）。取
%   $v=z\partial_z-w\partial_w$，则 $v(zw)=z\cdot w-w\cdot z=0$，成立；若取
%   $z\partial_z+w\partial_w$ 则得 $2zw\ne0$，不成立。故 $-w\partial_w$ 前的负号
%   是必需的（这是最容易抄错符号的地方）。]

由正合序列 (1.58)，我们得到下述推论。

\begin{corollary}[1.34]
下述 $\mathcal{O}_{\mathcal{C}}$-模序列
\[
0\to\Theta_{\mathcal{C}/\mathcal{B}}\to\Theta_{\mathcal{C}}
\to\pi^{-1}\Theta_{\mathcal{B}}\otimes\mathcal{O}_{\mathcal{C}}
\to\underline{\mathrm{Ext}}^1_{\mathcal{O}_{\mathcal{C}}}
(\Omega^1_{\mathcal{C}/\mathcal{B}},\mathcal{O}_{\mathcal{C}})\to 0
\]
是正合的。
\end{corollary}

% [原文 p.33]

\begin{lemma}[1.35]
存在正合序列
\begin{equation}
0\to\Theta_{\mathcal{B}}(-\log D)\to\Theta_{\mathcal{B}}
\xrightarrow{\ t\ }\pi_*\underline{\mathrm{Ext}}^1_{\mathcal{O}_{\mathcal{C}}}
(\Omega^1_{\mathcal{C}/\mathcal{B}},\mathcal{O}_{\mathcal{C}})\to 0,
\tag{1.62}
\end{equation}
其中
\[
\Theta_{\mathcal{B}}(-\log D)=\{v\in\Theta_{\mathcal{B}}\mid
v(\mathcal{I}_D)\subset\mathcal{I}_D\},
\]
而 $\mathcal{I}_D$ 是 $\mathcal{B}$ 中 $D$ 的定义理想层。
\end{lemma}
% [Tutor Note: $\Theta_{\mathcal{B}}(-\log D)$ 是"沿 $D$ 的对数向量场"层，即与
%   $D$ 相切的向量场；$D\subset\mathcal{B}$ 是判别除子（$C_s$ 奇异的 $s$ 的轨迹）。
%   局部核对：若 $D$ 由 $u=0$ 定义，则 $\mathcal{I}_D=(u)$，条件
%   $v(u)\in(u)$ 恰好排除 $\partial_u$ 方向而允许 $u\partial_u$，故
%   $\Theta_{\mathcal{B}}(-\log D)$ 局部由 $u\partial_u,\partial_{u_2},\dots$ 生成，
%   是秩 $\dim\mathcal{B}$ 的局部自由层。]

\noindent\textit{证明.}\quad 首先注意，层 $\Theta_{\mathcal{B}}(-\log D)$ 是
$\mathcal{B}$ 上与 $D$ 相切的向量场芽层。由于
$\pi:\mathcal{C}\to\mathcal{B}$ 是通用族，利用小平–Spencer 映射与 (1.53)，
对每一点 $s\in\mathcal{B}$ 我们有正合序列
\[
0\to H^1\Big(C_s,\Theta_{C_s}\Big(-\sum_{j=1}^{N}s_j(s)\Big)\Big)
\to T_s\mathcal{B}\to H^0\big(C_s,
\underline{\mathrm{Ext}}^1_{\mathcal{O}_{C_s}}
(\Omega^1_{C_s},\mathcal{O}_{C_s})\big)\to 0.
\]
$H^1\big(C_s,\Theta_{C_s}(-\sum_{j=1}^{N}s_j(s))\big)$ 的每个元素对应于
$\mathcal{B}$ 在 $s$ 处的一个保持 $C_s$ 的奇点的切向量。因此上述序列的层化
版本也是正合的。\hfill$\square$

\begin{theorem}[1.36]
设 $(\pi:\mathcal{C}\to\mathcal{B};s_1,s_2,\dots,s_N)$ 是亏格为 $g$ 的稳定
$N$-点曲线的通用族。则存在 $\mathcal{O}_{\mathcal{B}}$-模同构
\begin{equation}
\rho:\Theta_s(-\log D)\xrightarrow{\ \sim\ }
R^1\pi_*\big(\Theta_{\mathcal{C}/\mathcal{B}}(-S)\big), \tag{1.63}
\end{equation}
其中我们令 $S_j=s_j(\mathcal{B})$，$S=\sum_{j=1}^{N}S_j$。
\end{theorem}
% [Tutor Note: (1.63) 左端原文印作 $\Theta_s(-\log D)$，应为
%   $\Theta_{\mathcal{B}}(-\log D)$（证明的末句明确说 $\pi_*\mathcal{M}\cong
%   \Theta_{\mathcal{B}}(-\log D)$；且 $D\subset\mathcal{B}$，$s\notin\mathcal{B}$）。
%   照原文排，见报告 ERRATA。又 $S_j=s_j(\mathcal{B})$ 是 $\mathcal{C}$ 中的除子
%   （截面像），故 $\Theta_{\mathcal{C}/\mathcal{B}}(-S)$ 是在 $S$ 上消失的相对
%   向量场层，它才是全纯切丛上的层，右端才有 $R^1\pi_*$ 可言。]

\noindent\textit{证明.}\quad 将
$\underline{\mathrm{Hom}}_{\mathcal{O}_{\mathcal{C}}}(\,\cdot\,,
\mathcal{O}_{\mathcal{C}})$ 作用于正合序列 (1.58)，我们得到正合序列
\[
0\to\Theta_{\mathcal{C}/\mathcal{B}}\to\Theta_{\mathcal{C}}
\to\pi^{-1}\Theta_{\mathcal{B}}\otimes\mathcal{O}_{\mathcal{C}}
\to\underline{\mathrm{Ext}}^1_{\mathcal{O}_{\mathcal{C}}}
(\Omega^1_{\mathcal{C}/\mathcal{B}},\mathcal{O}_{\mathcal{C}})\to 0.
\]
该正合序列分裂成如下短正合序列：
\begin{equation}
0\to\Theta_{\mathcal{C}/\mathcal{B}}\to\Theta_{\mathcal{C}}
\xrightarrow{\ \kappa\ }\mathcal{M}\to 0, \tag{1.64}
\end{equation}
\begin{equation}
0\to\mathcal{M}\to\pi^{-1}\Theta_{\mathcal{B}}\otimes\mathcal{O}_{\mathcal{C}}
\to\underline{\mathrm{Ext}}^1_{\mathcal{O}_{\mathcal{C}}}
(\Omega^1_{\mathcal{C}/\mathcal{B}},\mathcal{O}_{\mathcal{C}})\to 0. \tag{1.65}
\end{equation}

设 $\mathcal{T}$ 为 $\mathcal{C}$ 上保持 $n$ 阶无穷小邻域的全纯向量场芽层。
层 $\mathcal{T}$ 由
\begin{equation}
\mathcal{T}=\{v\in\Theta_{\mathcal{C}}\mid
v(\mathcal{I}_S)\subset\mathcal{I}_S\} \tag{1.66}
\end{equation}
给出。
% [Tutor Note: (1.66) 中的 $\mathcal{I}_S$ 是除子 $S$ 的定义理想层。局部核对
%   （$S$ 由 $z=0$ 定义，$\mathcal{I}_S=(z)$）：条件 $v(z)\in(z)$ 允许
%   $v=z\partial_z$（$v(z)=z$）与 $v=\partial_{u_j}$（$v(z)=0$），排除
%   $v=\partial_z$（$v(z)=1\notin(z)$）。所以 (1.66) 与下文"$\mathcal{T}$ 由
%   $z\partial_z,\partial_{u_1},\dots,\partial_{u_M}$ 生成"完全一致。]

层 $\mathcal{T}$ 是 $\Theta_{\mathcal{C}}$ 的
$\mathcal{O}_{\mathcal{C}}$-子模，并且在
$\bigcup_{j=1}^{N}S_j$ 之外与 $\Theta_{\mathcal{C}}$ 重合。对点
$P\in S_j$，设 $(u_1,u_2,\dots,u_M,z)$ 为 $\mathcal{C}$ 的以 $P$ 为中心的
局部坐标，使得 $(u_1,u_2,\dots,u_M)$ 是 $\mathcal{B}$ 的以 $\pi(P)$ 为中心的
坐标，且 $S_j$ 在 $P$ 的邻域内由方程 $z=0$ 定义。于是，在 $P$ 的邻域内层
$\mathcal{T}$ 作为 $\mathcal{O}_{\mathcal{C}}$-模由
\[
z\frac{\partial}{\partial z},\quad
\frac{\partial}{\partial u_1},\ \dots,\
\frac{\partial}{\partial u_M}
\]
生成。因此 $\mathcal{T}$ 在 $\mathcal{C}$ 上是局部自由的。

让我们考察正合序列 (1.64) 与 (1.65)。由于层
$\underline{\mathrm{Ext}}^1_{\mathcal{O}_{\mathcal{C}}}
(\Omega^1_{\mathcal{C}/\mathcal{B}},\mathcal{O}_{\mathcal{C}})$ 的支集在
$\Sigma$ 中，层 $\mathcal{M}$ 在 $\mathcal{C}\setminus\Sigma$ 上等于
$\pi^{-1}\Theta_{\mathcal{B}}\otimes\mathcal{O}_{\mathcal{C}}$。利用
$\mathcal{C}$ 的以 $P\in S_j$ 为中心的上述局部坐标，(1.64) 中映射 $\kappa$
在 $P$ 的邻域内对 $\mathcal{T}$ 的限制由
\[
a(u,z)z\frac{\partial}{\partial z}
+\sum_j B_j(u,z)\frac{\partial}{\partial u_j}
\longmapsto\sum_j B_j(u,z)\frac{\partial}{\partial u_j}
\]
给出。
% [Tutor Note: 核对 $\kappa$ 的核：被映到 $0$ 的正是 $a(u,z)z\partial_z$，
%   即 $\mathcal{O}_{\mathcal{C}}\cdot z\partial_z\cong
%   \Theta_{\mathcal{C}/\mathcal{B}}(-S)$（因为 $\Theta_{\mathcal{C}/\mathcal{B}}$
%   在结点附近由 $z\partial_z-w\partial_w$ 生成，乘以 $(z)$ 后 $w\partial_w$ 项
%   落在 $(zw)=0$ 里，故只剩 $z\partial_z$）。这一步解释了下一行"核为
%   $\Theta_{\mathcal{C}/\mathcal{B}}(-(n+1)S)$"在 $n=0$ 时的情形。]

% [原文 p.34]

因此 $\kappa:\mathcal{T}\to\mathcal{M}$ 是满射，且其核在 $P$ 的邻域内为
$\Theta_{\mathcal{C}/\mathcal{B}}(-(n+1)S)$。另一方面，在
$\mathcal{B}\setminus\bigcup_{j=1}^{N}S_j$ 上，层 $\mathcal{T}$ 等于
$\Theta_{\mathcal{C}}$。于是我们得到正合序列
\begin{equation}
0\to\Theta_{\mathcal{C}/\mathcal{B}}(-S)\to\mathcal{T}
\to\mathcal{M}\to 0. \tag{1.67}
\end{equation}
% [Tutor Note: 原文此处有一个不一致：紧邻的文字说 $\kappa$ 的核是
%   $\Theta_{\mathcal{C}/\mathcal{B}}(-(n+1)S)$，而 (1.67) 写的是
%   $\Theta_{\mathcal{C}/\mathcal{B}}(-S)$。二者仅在 $n=0$ 时相同（$n$ 是
%   无穷小邻域的阶数，见 1.3 节）。照原文排，此处按 $n=0$ 理解 (1.67)，并把
%   一般 $n$ 的情形记为待核（见报告 ISSUES）。]

由正合序列 (1.67)，我们得到长正合序列
\begin{align}
0&\to\pi_*\big(\Theta_{\mathcal{C}/\mathcal{B}}(-S)\big)
\xrightarrow{\ \tau\ }\pi_*\mathcal{T}\to\pi_*\mathcal{M}
\xrightarrow{\ \rho\ }R^1\pi_*\big(\Theta_{\mathcal{C}/\mathcal{B}}(-S)\big)
\tag{1.68}\\
&\qquad\to R^1\pi_*\mathcal{T}\to R^1\pi_*\mathcal{M}\to 0.
\nonumber
\end{align}

令 $\mathcal{B}_0=\mathcal{B}\setminus D$，
$\mathcal{C}_0=\pi^{-1}(\mathcal{B}_0)$，$\pi_0=\pi|_{\mathcal{C}_0}$。
则在 $\mathcal{B}_0$ 上 $\pi_{0*}\mathcal{M}=\Theta_{\mathcal{B}}$，且由推论
1.26，同态 $\rho$ 就是小平–Spencer 映射。由于我们的族是通用族，$\rho$ 在
$\mathcal{B}_0$ 上是同构。因此，(1.68) 中的层同态 $\tau$ 在 $\mathcal{B}_0$
上是同构。但在 $\mathcal{B}_0$ 上我们有
$\pi_*(\Theta_{\mathcal{C}/\mathcal{B}}(-S))=0$。因此，$\pi_*\mathcal{T}$
在 $\mathcal{B}_0$ 上为 $0$。由于 $\mathcal{T}$ 局部自由，$\pi_*\mathcal{T}$
无挠，故在 $\mathcal{B}$ 上 $\pi_*\mathcal{T}=0$。这也推出在 $\mathcal{B}$ 上
\begin{equation}
\pi_*\big(\Theta_{\mathcal{C}/\mathcal{B}}(-S)\big)=0. \tag{1.69}
\end{equation}

下面我们证明 (1.68) 中的 $\rho$ 是同构。为此只需证明 $R^1\pi_*\mathcal{T}$
是局部自由的。因为若 $R^1\pi_*\mathcal{T}$ 局部自由，则由 $\rho$ 在
$\mathcal{B}_0$ 上是同构知 $\mathrm{Coker}\,\rho$ 是 $R^1\pi_*\mathcal{T}$
的挠子层，从而为零。由凝聚层的上同调理论，
\[
\chi(\mathcal{T}\otimes\mathcal{O}_{C_s})
=\dim_{\C}H^0(C_s,\mathcal{T}\otimes\mathcal{O}_{C_s})
-\dim_{\C}H^1(C_s,\mathcal{T}\otimes\mathcal{O}_{C_s})
\]
与 $s\in\mathcal{B}$ 无关，其中 $C_s=\pi^{(n)^{-1}}(s)$。（例如见 [BS]。）
此外，若 $\dim_{\C}H^1(C_s,\mathcal{T}\otimes\mathcal{O}_{C_s})$ 与 $s$
无关，记为 $k$，则由 $H^2(C_s,\mathcal{T}\otimes\mathcal{O}_{C_s})=0$ 知
$R^1\pi_*\mathcal{T}$ 是 $\mathcal{B}$ 上秩为 $k$ 的局部自由
$\mathcal{O}_{\mathcal{B}}$-模。因此，只需证明对一切 $s\in\mathcal{B}$ 都有
$H^0(C_s,\mathcal{T}\otimes\mathcal{O}_{C_s})=0$。
% [Tutor Note: 上同调维数跳变的判据（Mumford 的"半连续定理"）：若
%   $H^1$ 的维数在族中是常数，则 $R^1\pi_*$ 局部自由且与纤维上同调交换。
%   这里 $H^2=0$ 是因为 $\dim C_s=1$，$H^2$ 只能来自二维支撑（不存在）。
%   另外 $\chi$ 与 $s$ 无关用的是 Riemann–Roch（对局部自由层在纤维上的
%   Euler 示性数），所以先用 $\chi$ 定住 $H^0-H^1$，再用 $H^1$ 常数定住 $H^0$。]

由于 $C_s$ 是 $\mathcal{C}$ 中的局部完全交，我们有正合序列
\[
0\to\Theta_{C_s}\to\Theta_{\mathcal{C}}\otimes\mathcal{O}_{C_s}
\to\mathcal{O}_{C_s}(N)\to\underline{\mathrm{Ext}}^1_{\mathcal{O}_{C_s}}
(\Omega_{C_s},\mathcal{O}_{C_s})\to 0,
\]
其中 $N$ 是 $C_s$ 在 $\mathcal{C}$ 中的法丛，它是秩为 $3g-3+N$ 的平凡丛。
（例如见 [Ar]。）由此正合序列我们得到两个短正合序列
\[
0\to\Theta_{C_s}\to\Theta_{\mathcal{C}}\otimes\mathcal{O}_{C_s}
\to M_s\to 0,
\]
\[
0\to M_s\to\mathcal{O}_{C_s}(N)\to\underline{\mathrm{Ext}}^1_{\mathcal{O}_{C_s}}
(\Omega_{C_s},\mathcal{O}_{C_s})\to 0.
\]
与上面类似，我们有正合序列
\[
0\to\Theta_{C_s}\Big(-\sum_{j=1}^{N}Q_j\Big)
\to\mathcal{T}\otimes\mathcal{O}_{C_s}\to M_s\to 0,
\]
其中 $Q_j=s_j(s)$。这给出长正合序列
\begin{align}
0&=H^0\Big(C_s,\Theta_{C_s}\Big(-\sum Q_j\Big)\Big)
\to H^0(C_s,\mathcal{T}\otimes\mathcal{O}_{C_s})\to H^0(C_s,M_s)
\tag{1.70}\\
&\qquad\xrightarrow{\ \rho\ }H^1\Big(C_s,\Theta_{C_s}
\Big(-\sum Q_j\Big)\Big)\to H^1(C_s,\mathcal{T}\otimes\mathcal{O}_{C_s}).
\nonumber
\end{align}
% [Tutor Note: (1.70) 左端写 $0=H^0(\dots)$，是因为稳定曲线没有非平凡
%   无穷小自同构：$\deg\Theta_{C_s}(-\sum Q_j)=2-2g-N<0$，故 $H^0=0$。
%   这个"$H^0=0$"正是稳定性的解析表现，也是第三章讲共形块有限维性的基础。
%   注意 $\mathcal{O}_{C_s}(N)$ 记号表示法丛 $N$ 的截面层（不是线丛的
%   除子记号），因为 $N$ 平凡故同构于 $\mathcal{O}_{C_s}^{\oplus(3g-3+N)}$。]

上同调群 $H^0(C_s,M_s)$ 参数化 $C_s$ 在 $\mathcal{C}$ 中的无穷小位移。
（细节见 Tsuboi [Ts]。）由于 $\pi:\mathcal{C}\to\mathcal{B}$ 是通用族，
% [原文 p.35]

$C_s$ 在 $\mathcal{C}$ 中的无穷小位移与 $C_s$ 的无穷小形变相重合。因此，
(1.70) 中的同态 $\rho$ 是同构。于是我们有
\[
H^0(C_s,\mathcal{T}\otimes\mathcal{O}_{C_s})=0.
\]

最后我们证明 $\pi_*\mathcal{M}$ 同构于
$\Theta_{\mathcal{B}}(-\log D)$。由 (1.65) 我们得到正合序列
\[
0\to\pi_*\mathcal{M}\to\Theta_{\mathcal{B}}
\xrightarrow{\ t\ }\pi_*\big(\underline{\mathrm{Ext}}^1_{\mathcal{O}_{\mathcal{C}}}
(\Omega^1_{\mathcal{C}/\mathcal{B}},\mathcal{O}_{\mathcal{C}})\big).
\]
同态 $t$ 与正合序列 (1.62) 中出现的同态相同。故 $t$ 是满射。因此由引理
1.35，$\pi_*\mathcal{M}$ 同构于 $\Theta_{\mathcal{B}}(-\log D)$。
\hfill$\square$

\begin{remark}[1.37]
上述定理 1.36 中的同态 $\rho$ 也称为小平–Spencer 映射。上述证明表明存在
正合序列
\[
0\to\Theta_{\mathcal{C}/\mathcal{B}}\Big(-\sum S_j\Big)\to\mathcal{T}
\to\pi^{-1}\Theta_{\mathcal{B}}\otimes\mathcal{O}_{\mathcal{C}}
\to\underline{\mathrm{Ext}}^1_{\mathcal{O}_{C_s}}
(\Omega^1_{C_s},\mathcal{O}_{C_s})\to 0,
\]
其中 $\mathcal{T}$ 是 (1.66) 中定义的 $\Theta_{\mathcal{C}}$ 的子层。
% [Tutor Note: 这个序列右端的 $\underline{\mathrm{Ext}}^1_{\mathcal{O}_{C_s}}
%   (\Omega^1_{C_s},\mathcal{O}_{C_s})$ 按上文 (1.65) 应为
%   $\underline{\mathrm{Ext}}^1_{\mathcal{O}_{\mathcal{C}}}
%   (\Omega^1_{\mathcal{C}/\mathcal{B}},\mathcal{O}_{\mathcal{C}})$（在
%   $\mathcal{C}$ 上取，再在纤维上取 $H^0$）；原文把下标写成 $C_s$ 是记号滑落，
%   照原文排，见报告 ISSUES。]

选取一个小的 Stein 开集 $\mathcal{U}\subset\mathcal{B}$ 与向量场
$v\in H^0(\mathcal{U},\Theta_{\mathcal{B}}(-\log D))$。再选取
$\pi^{-1}(\mathcal{U})$ 的一个 Stein 开覆盖
$\{\mathcal{U}_j\}_{j\in J}$。于是 $v$ 也定义了一个元素
\[
\pi^*v\in H^0(\mathcal{U}_j,\pi^{-1}\Theta_{\mathcal{B}}
\otimes\mathcal{O}_{\mathcal{C}}),
\]
它在 $\underline{\mathrm{Ext}}^1_{\mathcal{O}_{C_s}}(\Omega^1_{C_s},
\mathcal{O}_{C_s})$ 中的像为零，因为切向量 $v$ 是保持奇点的无穷小形变的
方向。因此，若 $\mathcal{U}_j$ 足够小，我们可以找到元素
$v_j\in H^0(\mathcal{U}_j,\mathcal{T})$，它被映到 $\pi^*v$。于是我们有
\[
v_{ij}=v_j-v_i\in H^0\big(\mathcal{U}_i\cap\mathcal{U}_j,
\Theta_{\mathcal{C}/\mathcal{B}}(-S)\big),
\]
而 $\{v_{ij}\}$ 定义了一个元素
\[
[\{v_{ij}\}]\in H^1\big(\pi^{-1}(\mathcal{U}),
\Theta_{\mathcal{C}/\mathcal{B}}(-S)\big).
\]
映射
\[
v\longmapsto[\{v_{ij}\}]
\]
不是别的，正是定理 1.36 中的小平–Spencer 映射 $\rho$。
\end{remark}

% [Tutor Note: 这段给出了定理 1.36 的显式 Čech 描述：$\rho(v)$ 由
%   "把 $v$ 局部提升到 $\mathcal{T}$ 所得的转移函数"给出。核对
%   $v_{ij}\in\Theta_{\mathcal{C}/\mathcal{B}}(-S)$：$v_i,v_j$ 都提升同一个
%   $\pi^*v$，其差别落在 $\mathcal{T}$ 的核中，而由 (1.67)（$n=0$）核是
%   $\Theta_{\mathcal{C}/\mathcal{B}}(-S)$。这条显式公式在第五章算 Hitchin
%   联络的曲率时会重新用到。]

\begin{lemma}[1.38]
设 $\mathfrak{F}=(\pi:\mathcal{C}\to\mathcal{B};s_1,s_2,\dots,s_N)$ 是
亏格为 $g$ 的 $N$-点光滑曲线的通用族。设 $X,Y$ 是 $\mathcal{B}$ 的全纯
向量场，而
\[
\rho:\Theta_{\mathcal{B}}\to R^1\pi_*\Theta_{\mathcal{C}/\mathcal{B}}(-S)
\]
是小平–Spencer 映射。则我们有
\[
\rho([X,Y])=[\rho(X),\rho(Y)]+X(\rho(Y))-Y(\rho(X)),
\]
其中记
\[
\rho(X)=\{\theta_{\lambda\mu}\},\qquad
\rho(Y)=\{\tau_{\lambda\mu}\},
\]
则 $\rho([X,Y])$ 是由
\[
\vartheta_{\lambda\mu}=[\theta_{\lambda\mu},\tau_{\lambda\mu}]
\]
定义的上闭链 $\{\vartheta_{\lambda\mu}\}$，而 $X(\rho(Y))$ 是由
\[
\tau'_{\lambda\mu}=X(\tau_{\lambda\mu})
\]
定义的上闭链 $\{\tau'_{\lambda\mu}\}$。
\end{lemma}
% [Tutor Note: 公式 $\rho([X,Y])=[\rho(X),\rho(Y)]+X(\rho(Y))-Y(\rho(X))$
%   是"$\rho$ 不是 Lie 代数同态，而是一个'平坦性'关系"的精确表述。若把它
%   改写成 $\rho([X,Y])-X(\rho(Y))+Y(\rho(X))=[\rho(X),\rho(Y)]$，就认出
%   左边是某个联络的曲率零条件（$\nabla_X\nabla_Y-\nabla_Y\nabla_X
%   -\nabla_{[X,Y]}$ 的一阶项）；第五章证明 Hitchin 联络射影平坦时，
%   这个公式是零阶部分的来源。注意 $\rho(X)$ 是一次上闭链（Čech 1-上链），
%   而 $X(\rho(Y))$ 是对上链的系数逐项求导，仍在同一层里。]

% [原文 p.36]

\noindent\textit{证明.}\quad 我们可以假定 $\mathcal{B}$ 足够小。设 $\{U_{\lambda\mu}\}$
是 $\mathcal{C}$ 的一个足够细的开覆盖。设 $(z_1,\dots,z_m)$ 是
$\mathcal{B}$ 的坐标，而 $(z_1,\dots,z_m,\xi_\lambda)$ 是 $U_\lambda$ 的
局部坐标。此外，我们可以假定：对 $j\in\{1,2,\dots,N\}\subset\Lambda$，
$U_j$ 是 $s_j(\mathcal{B})$ 的坐标邻域，使得 $\xi_j$ 诱导出 $n$ 阶无穷小
邻域 $\eta_j$，且 $s_j(\mathcal{B})$ 由 $\xi_j=0$ 定义。对
$U_\lambda\cap U_\mu\ne\emptyset$，我们有
\[
\xi_\lambda=h_{\lambda\mu}(\xi_\mu,z_1,\dots,z_m).
\]

向量场 $X$ 诱导坐标 $(z_1,\dots,z_m)$ 的一个无穷小变换：
\[
(z_1,\dots,z_m)\longmapsto
(z_1+\epsilon_1a_1(z),\dots,z_m+\epsilon_1a_m(z)),
\]
其中 $\epsilon_1$ 是一个对偶数，即 $\epsilon_1^2=0$，且
\[
X=\sum_{j=1}^{m}a_i(z)\frac{\partial}{\partial z_i}.
\]
% [Tutor Note: 求和指标写成 $j$ 而下标是 $i$ 是原文的下标笔误，应为
%   $X=\sum_{j=1}^{m}a_j(z)\partial/\partial z_j$（否则 $i$ 无定义）。
%   照原文排，见报告 ERRATA。对偶数 $\epsilon$ 满足 $\epsilon^2=0$，
%   用它做"无穷小形变"就等价于在 $\C[\epsilon]/(\epsilon^2)$ 上工作，
%   这避免了 $\epsilon\to0$ 的极限语言。]

为简单起见，我们把 (1.4) 写成 $z\mapsto z+\epsilon_1X$。由于我们的族是
通用族，我们有
\[
g_{\lambda\mu}(\xi_\mu,z+\epsilon_1X)
=g_{\lambda\mu}(\xi_\mu,z)+\epsilon_1\ell_{\lambda\mu}(\xi_\mu,z),
\]
其中我们令
\[
\theta_{\lambda\mu}=\ell_{\lambda\mu}\frac{\dd}{\dd\xi_\mu},\qquad
\tau_{\lambda\mu}=m_{\lambda\mu}\frac{\dd}{\dd\xi_\mu}.
\]
% [Tutor Note: $g_{\lambda\mu}$ 是 $U_\lambda$ 与 $U_\mu$ 之间的转移函数
%   （坐标变换），$\ell_{\lambda\mu},m_{\lambda\mu}$ 分别是沿 $X$、$Y$ 方向
%   形变时转移函数的一阶变分。$\theta,\tau$ 是"相对向量场值"的
%   Čech 1-上链，取值在 $\Theta_{\mathcal{C}/\mathcal{B}}$ 中；这里把
%   $\dd/\dd\xi_\mu$ 作为纤维方向的基（在 $s_j$ 附近 $\xi$ 就是 $z$）。
%   原文 $\theta$ 用 $\dd/\dd\xi_\mu$、后面 (p.36 末) 又用
%   $\dd/\dd\xi_\lambda$，是同一对象的两种写法。]

类似地，我们有
\[
g_{\lambda\mu}(\xi_\mu,z+\epsilon_2Y)
=g_{\lambda\mu}(\xi_\mu,z)+\epsilon_2m_{\lambda\mu}(\xi_\mu,z),
\]
其中 $\epsilon_2$ 是另一个对偶数。现在让我们计算先沿 $X$ 的方向作无穷小
形变、再沿 $Y$ 的方向作无穷小形变所得的结果：
\begin{align*}
&g_{\lambda\mu}(\xi_\mu,z+\epsilon_2Y)
+\epsilon_1\ell_{\lambda\mu}(\xi_\lambda+\epsilon_2m_{\lambda\mu},
z+\epsilon_2Y)\\
&\quad=g_{\lambda\mu}(\xi_\mu,z)+\epsilon_2m_{\lambda\mu}(\xi_\lambda,z)
+\epsilon_1\ell_{\lambda\mu}(\xi_\lambda,z)
+\epsilon_1\epsilon_2\Big(Y(\ell_{\lambda\mu}(\xi_\lambda,z))
+m_{\lambda\mu}\frac{\dd\ell_{\lambda\mu}}{\dd\xi_\lambda}\Big).
\end{align*}
若按相反的顺序形变，我们有
\begin{align*}
&g_{\lambda\mu}(\xi_\mu,z)+\epsilon_1\ell_{\lambda\mu}(\xi_\lambda,z)
+\epsilon_2m_{\lambda\mu}(\xi_\lambda,z)\\
&\quad+\epsilon_1\epsilon_2\Big(X(m_{\lambda\mu}(\xi_\lambda,z))
+\ell_{\lambda\mu}\frac{\dd\ell_{\lambda\mu}}{\dd\xi_\lambda}\Big).
\end{align*}

因此，若令
\[
\rho([X,Y])=\{\tilde{\vartheta}_{\lambda\mu}\},
\]
我们有
\begin{align*}
\tilde{\vartheta}_{\lambda\mu}
&=\Big\{\ell_{\lambda\mu}\frac{\dd m_{\lambda\mu}}{\dd\xi_\lambda}
-m_{\lambda\mu}\frac{\dd\ell_{\lambda\mu}}{\dd\xi_\lambda}
+X(m_{\lambda\mu})-Y(\ell_{\lambda\mu})\Big\}
\frac{\dd}{\dd\xi_\lambda}\\
&=[\theta_{\lambda\mu},\tau_{\lambda\mu}]
+\big\{X(\tau_{\lambda\mu})-Y(\ell_{\lambda\mu})\big\}
\frac{\dd}{\dd\xi_\lambda}\\
&=[\rho(X),\rho(Y)]+X(\rho(Y))-Y(\rho(X)).
\end{align*}
\hfill$\square$
% [Tutor Note: 中间一行原文印作 $X(\tau_{\lambda\mu})-Y(\ell_{\lambda\mu})$，
%   其中第二项应为 $Y(\theta_{\lambda\mu})$（或写成
%   $\{X(m_{\lambda\mu})-Y(\ell_{\lambda\mu})\}\dd/\dd\xi_\lambda$）：因为
%   第一行里的一次项就是 $X(m)-Y(\ell)$，而
%   $X(\tau_{\lambda\mu})=X(m_{\lambda\mu})\dd/\dd\xi_\lambda$，
%   $Y(\theta_{\lambda\mu})=Y(\ell_{\lambda\mu})\dd/\dd\xi_\lambda$。
%   照原文排，见报告 ERRATA。最后一行用到
%   $[\theta,\tau]=(\ell\,\dd m/\dd\xi-m\,\dd\ell/\dd\xi)\dd/\dd\xi$
%   （对偶数环上 $\dd/\dd\xi$ 与函数相乘的通常交换关系）。]
